% SPDX-License-Identifier: Apache-2.0
\section{What the Runtime Is}
\label{sec:r-overview}

The embedding article~\cite{embedding} argues that a hardware
description language embedded in Rust needs no parser and no keywords,
and that what remains is a library. The crate \code{txhdl} is that
library, and this document is the crate. Every file below is included
from the tree when the document is built, so a listing here cannot
differ from the code, and the rule the project keeps is that a concept
enters the runtime, a compiled example and these documents in the same
change, or the change is not finished.

The crate is small on purpose. Nine modules and one procedural macro
crate hold everything:

\begin{itemize}
\item \code{txhdl::types}: two-valued \code{Bit}, nine-valued
  \code{Logic}, the numeric vectors \code{U<N>} and \code{I<N>},
  \code{logic::Vec<N>}, and the two marker traits \code{Transaction}
  and \code{Tag}.
\item \code{txhdl::comp}: clocks, wires and channels and their ends,
  interfaces, crossings, registers and memories, units, configurations,
  and the executor.
\item \code{txhdl::pipeline}: operators that take a cycle, all
  \code{async}.
\item \code{txhdl::funcs}: operators that cannot, all plain \code{fn}.
\item \code{txhdl::netlist}: the lowering's target, an expression
  tree, statements, a \code{Lowered} unit and its two emitters,
  Verilog and VHDL; and the structural skeleton from the trace walk.
\item \code{txhdl::foreign}: a Verilog or VHDL module as a unit, run
  beside the Rust ones.
\item \code{txhdl::map}: an address map as a trait constant.
\item \code{txhdl::regmap}: what a map declared with \code{regmap!}
  stands for at run time.
\item \code{txhdl::formal}: \code{check!}, \code{assume!} and
  \code{cover!}.
\item \code{txhdl\_macros}: the five derives, \code{interface!},
  \code{with!}, \code{when!}, \code{case!}, \code{station!},
  \code{regmap!}, \code{\#[pipeline]} and \code{\#[lower]}, without \code{syn} or \code{quote}; and
  \code{select!}, a \code{macro\_rules!} in the crate root.
\end{itemize}

The crate is a simulation-shaped prototype: values exist while it
runs, wires are shared cells, and a testbench may look at anything.
The lowering to hardware is the experiment the article states, and
this crate is what it was written against; the experiment has been
run, Section~\ref{sec:r-lowering}.

\subsection{The model of time}
\label{sec:r-time}

The executor is the one part of the runtime a reader cannot recover
from the signatures, so it is stated here before the code.

\textbf{Time runs in ticks.} One poll of the top unit is one tick.
Every clock is a type with a period and a phase in ticks, and has a
rising edge at every tick \code{t} with \code{t >= PHASE} and
\code{(t - PHASE) \% PERIOD == 0}, and a falling edge half a period
later. The default clock has period 2 and phase 0, so its cycle is two
ticks and its falling edge is a tick of its own; a single-clock design
mentions neither number, and its testbench steps by
\code{Running::cycle}.

\textbf{A process waits, then reads.} The waits are events:
\code{C::rising()} and \code{C::falling()}, the next edge of a clock;
\code{rx.wait()}, the next rising edge of the channel's clock at which
the channel holds a transaction, taken; \code{tx.put(|| v)}, its
sender's side, the next such edge at which the channel has room, and
\code{v} sent in the step that edge begins, so a register set after it
takes its value on the edge the transaction is taken (issue 755); and
\code{until(C::rising, || cond)},
or the same with \code{falling}, the next such edge at which a
condition on state holds. Reads are plain: a register in an
expression, or \code{Reg::get}, is what the last edge latched, since
a register is \code{Copy} and has its value's operators; a wire's
\code{get} returns its driven value, a memory's \code{read} accesses
its asynchronous read port. Drives are
deferred: \code{Reg::set} and \code{Mem::write} take effect at the end
of the step, so a read after a drive in the same iteration still sees
what the edge latched. A memory refuses an address past its end, a
read when it is made and a write when it happens, as the netlist does:
VHDL indexes the array as written and nvc stops on an index out of
range, and a runtime that wrapped the address agreed with Verilator on
a word the hardware never had (issue 556).

\textbf{A memory a block RAM cannot hold is refused when it is
lowered.} A block RAM has two ports, and a read at the address a write
is at shares the write's port, as a read-modify-write does. A memory
of more than 4096 bits that is written at two addresses, or reached at
three, Vivado builds from LUTs, which no test sees and a synthesis an
hour later shows as lost timing. The lowering counts each memory's
read and write addresses across the unit's processes and wires, and
refuses such a memory with a panic that names the unit, the memory and
the addresses (issue 1285). The remedy is to read it at one address a
cycle beside the write's, through a multiplexer of the addresses, or
to mark the field \code{\#[distributed]} when LUT RAM is what is
meant. A memory of 4096 bits or fewer, a register file or a small
FIFO, is left alone.

\textbf{A memory can say what Vivado should make it.}
\code{\#[ram\_style("block")]} on a memory's field, or
\code{"distributed"}, \code{"registers"} or \code{"ultra"}, puts
the \code{ram\_style} attribute on the memory in both netlists:
\code{(* ram\_style = "block" *)} before its declaration in Verilog,
and the attribute on its signal in VHDL (issue 1371).
\code{\#[distributed]} says \code{"distributed"}. Vivado chooses a
memory's kind from its shape, and a memory whose read is registered
can still come out as LUT RAM, as Razboj's colour bank, 4096 words of
32 bits, does (issue 1343). The attribute is a request: what a block
RAM cannot hold Vivado still refuses, and says so in its synthesis
log as an infeasible \code{ram\_style}, which is where to look. A
memory that says nothing is left to Vivado, as before.

\textbf{A register can say whether Vivado may use DSP slices.}
\code{\#[use\_dsp("no")]} on a register's field, or \code{"yes"},
\code{"logic"} or \code{"simd"}, puts the \code{use\_dsp} attribute
on it in both netlists (issue 1383). It is what a block RAM's read
register beside a multiply needs: Vivado pulls up to two registers on
a path into a multiply's input, and the read's register is one of
them, which leaves the read asynchronous and the memory LUT RAM, even
where the multiply ends in LUTs (issue 1343). \code{"no"} on the
products keeps the read's register where it is, at no cycle's cost.

\textbf{A channel is an elastic buffer of two.} A send is an offer:
\code{valid} and \code{data} high this step, the transaction in the
buffer at the end of it; \code{ready} says whether there is room, as
the edge left it. A receive is a take: \code{ready} high this step,
the head of the buffer gone at the end of it; \code{peek} is the head
as the edge left it, and \code{take} is the take of whatever is
offered, whether one was and the value, for a process that serves
the channel every cycle. Both commit as a register drive does, so what a
process sees of a channel does not depend on which process ran first
in the step, and two units that both run every cycle pass one
transaction per cycle with \code{valid} and \code{ready} registered
on both sides and no combinational path between them. That is what a
lowered unit's channel port sees and drives, which is why a unit with
channels can be checked as a netlist against its trace. A channel's
trace is its six wires by side: \code{rx\_data}, \code{rx\_valid} and
\code{rx\_ready} on the receiver's, \code{tx\_data}, \code{tx\_valid}
and \code{tx\_ready} on the sender's.

\textbf{An unregistered channel passes an offer in the cycle it is
made.} \code{chan\_unregistered}, which \code{\#[unregistered]} on a
\code{let} of \code{chan} in a unit of units makes, is the same
buffer of two with one bypass (issue 1293). While the buffer is empty,
the receiver's \code{peek}, \code{head} and \code{recv} see the
sender's offer of this step, and a take of it takes it this step, so it
never enters the buffer. An offer the receiver does not take is
buffered at the edge, as a registered channel's is. \code{ready} stays
the buffer as the edge left it, so only \code{valid} and \code{data}
cross in the cycle, and nothing runs from the receiver back to the
sender. The step polls each process once, so the receiver must run
after its sender: a receiver that found the channel empty before its
sender offered in the same step is refused when the sender offers, with
a panic naming the line that made the channel, rather than see the
offer a step later than the netlist would. In the netlist the channel
is the second channel module, \code{<top>\_txhdl\_chan\_u}, whose
\code{rx\_valid} and \code{rx\_data} are the sender's while the buffer
is empty. A unit of units whose unregistered channels close a ring of
wires through its children is refused when it is lowered, naming the
ring; a child's paths are read from its own lowering, so a child that
answers from a register breaks a ring.

\textbf{A wire kept as a field is a \code{let} with a name in the
trace.} \code{Wire<T>} is driven with \code{set} in the step, may be
read back with \code{get} in the same step, and is set before it is
read: it keeps the last value set, so a \code{get} before the step's
\code{set} gives the previous step's value, which the netlist's wire
would not; a unit's trace shows it under the unit's name as it shows a
register, and the lowering declares it from the field and drives it
as the wire the \code{let} would have been, so an internal signal is
seen in a waveform, and checked in the netlist, with no port for it.

\textbf{Every \code{.await} is a cycle boundary.} A process that has
crossed an edge at this step cannot cross it again, so a second wait
written after a first waits for the next edge, as an operator does.
Waits that share one edge are written inside \code{parallel!}, which
polls its futures together and returns their outputs as a tuple; that
is the one case the executor treats specially. Once one wait of the
group has crossed the edge at this step, the others in the group are
free, each once, so the group is one wait. State read before a wait,
at the top of a loop, is read in the previous step and before that
step's drives commit; state is read after the wait.

\textbf{A process is its waker.} The executor tells processes apart by
the data pointer of the waker they are polled with. \code{join2} and
\code{join\_all} give each child a fresh one, so the rule above holds
per process rather than per unit. \code{parallel!} keeps its futures in
the one process, which is the difference between the two: \code{join2}
forks, \code{parallel!} does not.

\textbf{An operator's cycle is in the process's clock.} An operator in
\code{txhdl::pipeline} waits one edge of whatever clock the process
last crossed an edge of, which is the only sense a latency has inside
a process. A process that has crossed no edge yet is in the default
clock.

\textbf{What the executor does not check.} A process that reads a
register of another clock other than through \code{ChanCdc} is a
hazard, and the prototype runs it. The register's type says its clock;
the process has no type that says one. \code{ChanCdc}, which the
asynchronous FIFO example runs, commits this hazard on purpose, on its
memory, under the pointer discipline that makes it safe.

\subsection{What is prototype only}

Two things in the code exist for the simulation and would not be
lowered. \code{Running} and \code{simulate} drive a design from a
\code{main}, and \code{now} reports the time step for traces. And the
shared cells behind every wire, channel and register, with the deferred
drives that commit at the end of a step, are how a simulation shares
state; the lowering replaces them with nets and flops.

\subsection{Tracing}

A design can be watched as a waveform. \code{comp::trace} holds a
registry of probes, each a dotted path, a width and a closure that
reads the signal; \code{\#[derive(Trace)]} on a unit registers every
field under its field name, so the hierarchy is the nesting of units
and nothing is named at construction. Registers, wire ends, channel
ends and crossings register themselves, plain values register nothing,
and \code{\#[derive(Value)]} gives a struct or a fieldless enum of the
design's own its width and bits. \code{Vcd} is the sink: the executor
calls it at the end of every tick, after the drives have committed,
and it writes what changed; a clock is written as the level it is,
high from its rising edge to its falling one. \code{Fst} is the
same probes with another writer, over the \code{fst-writer} crate,
the project's first external dependency, pinned by lockfile through
\code{crate\_universe}. \code{Wave::from\_env} chooses the sink
by the environment, \code{TXHDL\_FST} or \code{TXHDL\_VCD}, and
\code{stop} finishes the file. A compound value is one bit vector and
one signal per field beside it, from \code{Value::parts}; the writer
has no string variables, so a value keeps its bits.
The examples document draws the FST through \code{vcdcvt} and
\code{sqlite2drawtiming}, prebuilt tools pinned by checksum, and
\code{//tools/dt2tikz}, which turns drawtiming's text into TikZ.

\subsection{Netlists}
\label{sec:r-lowering}

The probe registry records, for every signal, its kind and the shared
cell it is on, so the two ends of one wire match. \code{netlist::verilog}
walks a design through it and writes a Verilog skeleton: a module per
unit, a \code{reg} per register, ports where a unit holds one end of
a wire whose other end is elsewhere, and instances with their ports
connected. It sees fields only, so a unit whose ports are parameters
of \code{run} yields its registers and nothing else, and it writes no
behaviour. That is the structural half of the lowering. The
behavioural half is the rest: \code{\#[pipeline]} lowers an
\code{async fn} of awaited operators, and \code{\#[lower]} lowers a
unit's \code{run} loop into a \code{Lowered} unit built from the
unit's fields, the ports of \code{run} and the statements the macro
read: registers, memories, wires and channel ports; \code{with!},
\code{when!}, \code{case!}, \code{if} and \code{select!}; slices,
concatenation and
extension; a \code{let} as a wire named for it. The macro builds an
expression tree, \code{netlist::Expr}, and two emitters render the
unit, Verilog and VHDL-2008. Both close the loop: the build simulates
a lowered unit's VHDL under \code{nvc} and its Verilog under
Verilator, each with a testbench \code{//tools/fst2tb} generates from
the trace the Rust run wrote, and each test passes only if registers
and outputs agree at every tick.
A unit's processes are its loops, one or several under
\code{join2}, each on the rising or the falling edge of its clock,
and each a clocked block; memories lower; and a unit with channels
is checked as any other, since the runtime's channel is a buffer
with a registered handshake and the netlist's port is what that
buffer shows. What the macro does not read it refuses by name.

\subsection{Building it}

\begin{lstlisting}[style=ascii,caption={The runtime and everything written against it, built hermetically.},label={lst:r-build}]
bazel build //lib/...                 # the crate and the macros
bazel build //lib/examples/...        # every example
bazel run //lib/examples:ex_fifo      # one that prints a trace
bazel run @rules_rust//:rustfmt \
  --@rules_rust//:rustfmt.toml=//:rustfmt.toml -- //lib/...
\end{lstlisting}

The last line formats the tree at 80 columns, which is the width these
documents include a file at. A line longer than that overflows its
frame, visibly, rather than wrapping.
